\documentclass[11pt]{article}

\usepackage[T1]{fontenc}
\usepackage{amsmath,amssymb,amsthm,mathtools}
\usepackage[a4paper,margin=30mm]{geometry}
\usepackage[hidelinks]{hyperref}
\usepackage{microtype}

\newtheorem{theorem}{Theorem}[section]
\newtheorem{lemma}[theorem]{Lemma}
\newtheorem{remark}[theorem]{Remark}
\newcommand{\R}{\mathbb{R}}
\newcommand{\linf}{\ell_\infty}
\newcommand{\ip}[2]{\langle #1,#2\rangle}

\title{A Factor-Two Stability Bound for\\Maximal Relative Projection Constants}
\author{Xiaojian Zeng\\
  \small PBC School of Finance, Tsinghua University, 100190\\
  \small \texttt{xiaojianzeng@foxmail.com}}
\date{}

\begin{document}
\maketitle

\begin{abstract}
For integers $n\ge r\ge1$, let $\lambda(r,n)$ be the maximal relative
projection constant of an $r$-dimensional subspace of $\ell_\infty^n$, and
let $\lambda(r)$ be the maximal absolute projection constant in real
dimension $r$.  Kumar, Mohar, Mojallal, and Pragada recently proved that
$\lambda(r,n)=\lambda(r)$ whenever
$n\ge 2^r\binom{r+1}{2}$.  We prove the same conclusion under the
factor-two smaller threshold
\[
 n\ge 2^{r-1}\binom{r+1}{2}.
\]
The additional observation is projective: the rows in the maximizing
matrix formulation may be independently multiplied by signs without
changing feasibility or the objective.  Normalizing one coordinate
therefore leaves only half of the coordinate orthants.  We verify that the
normalization is compatible with the cloning argument, including rows on
orthant boundaries and zero rows.
\end{abstract}

\section{Statement and context}

All spaces below are real.  If $Y$ is a subspace of a normed space $X$, set
\[
 \lambda(Y,X)=\inf\{\lVert P\rVert:P:X\to Y\text{ is a projection}\}.
\]
For a finite-dimensional normed space $Y$, its absolute projection constant
is
\[
 \lambda(Y)=\sup\{\lambda(TY,X):T:Y\to X\text{ is a linear isometry}\}.
\]
Projection constants in this sense go back at least to
Grünbaum~\cite{Grunbaum1960}.  For integers $n\ge r\ge1$, define
\[
 \lambda(r,n)=\sup\{\lambda(Y,\ell_\infty^n):
       Y\subseteq\ell_\infty^n,\ \dim Y=r\},
 \qquad
 \lambda(r)=\sup_{\dim Y=r}\lambda(Y).
\]
The sequence $\lambda(r,n)$ is nondecreasing in $n$ and has supremum
$\lambda(r)$; see, for example, the discussion and references in
\cite{FoucartSkrzypek2017,KumarEtAl2026}.  Basso's proposed general
stabilization argument~\cite{Basso2019} was subsequently shown to be
incomplete in the corresponding corrigendum~\cite{Basso2024}.

Kumar, Mohar, Mojallal, and Pragada~\cite{KumarEtAl2026} proved stability
for $n\ge2^r\binom{r+1}{2}$ by combining a maximizing-pair formulation with
a cloning reduction over all $2^r$ coordinate orthants.  Their paper is the
closest prior result to the theorem below.  It notes that minor improvements
may be possible, but does not state the factor-two bound.

\begin{theorem}\label{thm:main}
For every integer $r\ge1$ and every integer
\[
 n\ge 2^{r-1}\binom{r+1}{2},
\]
one has $\lambda(r,n)=\lambda(r)$.
\end{theorem}

The proof exploits the fact that the optimization depends only on absolute
inner products: each row vector can be reoriented independently.  Sign
switching of Gram vectors also occurs in other sharp Gram-matrix problems,
for example in the equality classification for rank-two correlation
matrices in~\cite{ZengLiuRatnavelu2026}; that result concerns permanents and
does not imply a statement about projection constants.  Here the reorientation
halves an orthant count while preserving every part of the maximizing-pair
argument.

\section{Maximizing pairs and cloning}

We recall the finite optimization used in the preceding stability proof.
For $U\in\R^{n\times r}$, write $x_i^T$ for its $i$th row.  Let
$\mathcal Q_n$ be the set of symmetric $n\times n$ sign matrices with
diagonal entries equal to one, and let $D_w=\operatorname{diag}(w)$.
Set
\[
 \mathcal W_n=\{w\in\R_{\ge0}^n:\lVert w\rVert_2=1\},
 \qquad
 \mathcal U_{n,r}=\{U\in\R^{n\times r}:U^TU=I_r\}.
\]
The Chalmers--Lewicki characterization, in the form used in
\cite[Theorem~2.1]{KumarEtAl2026} (see also
\cite{ChalmersLewicki2010,FoucartSkrzypek2017}), gives the two equal maxima
\begin{align}
 \lambda(r,n)
 &=\max_{w\in\mathcal W_n,\ Q\in\mathcal Q_n}
      \sum_{k=1}^r\lambda_k(D_wQD_w) \label{eq:eigen-optimization}\\
 &=\max_{U\in\mathcal U_{n,r},\ w\in\mathcal W_n}
   \sum_{i,j=1}^n w_iw_j\lvert\ip{x_i}{x_j}\rvert .
   \label{eq:optimization}
\end{align}
Call a feasible pair $(U,w)$ \emph{maximizing} if it attains this maximum,
and denote the displayed objective by $\sigma(U,w)$.

For completeness, we record the precise cloning fact needed later.  This is
the content of the cloning argument in~\cite[Claims~2.1--2.4]{KumarEtAl2026},
with the preservation details included.

\begin{lemma}[Cloning]\label{lem:cloning}
Let $(U,w)$ be a maximizing pair in \eqref{eq:optimization}.  Suppose
$C\subseteq\{1,\ldots,n\}$ satisfies
\[
 \ip{x_i}{x_j}\ge0\quad(i,j\in C),
\]
and suppose $c\in\R_{\ge0}^n$ satisfies $c_i=1$ for $i\notin C$ and
\begin{equation}\label{eq:moment-preserve}
 \sum_{i\in C}c_i x_ix_i^T=\sum_{i\in C}x_ix_i^T.
\end{equation}
Then the rows $\widetilde x_i^T=\sqrt{c_i}\,x_i^T$ and weights
$\widetilde w_i=\sqrt{c_i}\,w_i$ form another maximizing pair.
\end{lemma}

\begin{proof}
Let $B=(\lvert\ip{x_i}{x_j}\rvert)_{i,j}$.  Maximization first in $w$,
together with Perron--Frobenius, gives
\begin{equation}\label{eq:B-eigenvector}
 Bw=\lambda(r,n)w,
 \qquad \lambda_1(B)=\lambda(r,n).
\end{equation}
Choose the sign matrix $Q$ by $Q_{ij}=1$ if
$\ip{x_i}{x_j}\ge0$ and $Q_{ij}=-1$ otherwise, and put
$M=D_wQD_w$.  For this choice,
$\operatorname{tr}(U^TMU)=\sigma(U,w)=\lambda(r,n)$.  The Ky--Fan maximum
principle and \eqref{eq:eigen-optimization} show that both intervening
inequalities are equalities.  Hence the columns of $U$ span an invariant
subspace of $M$.  Thus there is a symmetric $H\in\R^{r\times r}$ with
$MU=UH$.  Taking the inner product of row $i$ of this identity with $x_i$,
and using \eqref{eq:B-eigenvector}, yields
\begin{equation}\label{eq:quadratic-identity}
 x_i^THx_i=\lambda(r,n)w_i^2\qquad(1\le i\le n).
\end{equation}

Equation \eqref{eq:moment-preserve} gives
$\widetilde U^T\widetilde U=U^TU=I_r$.  Taking its trace after
multiplication by $H$ and using \eqref{eq:quadratic-identity} gives
\[
 \lambda(r,n)(\lVert\widetilde w\rVert_2^2-\lVert w\rVert_2^2)
 =\operatorname{tr}\!\left(
 H\sum_{i\in C}(c_i-1)x_ix_i^T\right)=0.
\]
Since $\lambda(r,n)\ge\lambda(r,r)=1$, the new weights are again unit.

Finally, expanding the change in the objective gives
\begin{align*}
 \sigma(\widetilde U,\widetilde w)-\sigma(U,w)
 &=2\sum_{i\in C}(c_i-1)w_i
       \sum_{j=1}^n w_j\lvert\ip{x_i}{x_j}\rvert\\
 &\quad+\sum_{i,j\in C}(c_i-1)(c_j-1)w_iw_j
       \lvert\ip{x_i}{x_j}\rvert.
\end{align*}
The first term is zero by \eqref{eq:B-eigenvector} and the already proved
identity $\lVert\widetilde w\rVert_2=\lVert w\rVert_2$.  Because inner
products within $C$ are nonnegative, the second term equals
\[
 \left\lVert\sum_{i\in C}(c_i-1)w_ix_i\right\rVert_2^2\ge0.
\]
Thus $\sigma(\widetilde U,\widetilde w)\ge\lambda(r,n)$.  Feasibility and
\eqref{eq:optimization} force equality.
\end{proof}

\section{Half-orthant reduction}

Put
\[
 d=\dim\mathbb S^r=\binom{r+1}{2},
 \qquad N_r'=2^{r-1}d,
\]
where $\mathbb S^r$ is the space of real symmetric $r\times r$ matrices.
Fix any $n\ge r$ and a maximizing pair $(U,w)$.  For each row $x_i^T$ set
\[
 \varepsilon_i=\left\{
 \begin{array}{ll}
  1,&(x_i)_1\ge0,\\
  -1,&(x_i)_1<0,
 \end{array}\right.
 \qquad D=\operatorname{diag}(\varepsilon_1,\ldots,\varepsilon_n),
 \qquad V=DU.
\]
Writing $y_i^T=\varepsilon_i x_i^T$ for the rows of $V$, we have
\begin{equation}\label{eq:sign-invariance}
 V^TV=U^TD^2U=I_r,
 \qquad
 \lvert\ip{y_i}{y_j}\rvert=\lvert\ip{x_i}{x_j}\rvert.
\end{equation}
Therefore $(V,w)$ is still maximizing, and $(y_i)_1\ge0$ for all $i$.

Use the convention $\operatorname{sgn}(t)=1$ for $t\ge0$ and
$\operatorname{sgn}(t)=-1$ for $t<0$.  The normalized rows occupy only the
$2^{r-1}$ sign cells
\[
 \mathcal S_+=\{s\in\{\pm1\}^r:s_1=1\}.
\]
For $s\in\mathcal S_+$ let $C_s$ consist of the indices whose current rows
have sign vector $s$.  If $i,j\in C_s$, every coordinatewise product is
nonnegative; hence $\ip{y_i}{y_j}\ge0$, including when one or more
coordinates vanish.

Consider the moment matrix
\[
 A_s=\sum_{i\in C_s}y_iy_i^T\in\mathbb S^r.
\]
The conic form of Carathéodory's theorem (see, e.g.,
\cite[Corollary~7.1i]{Schrijver1986}) represents $A_s$ as
\begin{equation}\label{eq:caratheodory}
 A_s=\sum_{i\in I_s}\beta_i y_iy_i^T,
 \qquad I_s\subseteq C_s,\quad |I_s|\le d,\quad\beta_i\ge0.
\end{equation}
Zero generators may be omitted, and when $A_s=0$ we use the empty sum.
Set $c_i=\beta_i$ on $I_s$, $c_i=0$ on $C_s\setminus I_s$, and
$c_i=1$ off $C_s$.  Equations \eqref{eq:caratheodory} and
Lemma~\ref{lem:cloning} replace the current maximizing pair by one having at
most $d$ nonzero rows in this cell.

Repeat the operation for every $s\in\mathcal S_+$.  This iteration needs a
minor boundary check.  Multiplication by a positive scalar preserves the
complete sign vector of a nonzero row.  Multiplication by zero creates only
the zero row, which belongs to the all-positive cell under our convention
but remains zero under all later cloning operations.  Thus a later step
cannot create a nonzero row in an already processed cell.  Moreover, after
every step the pair is maximizing, so the graph, the invariant subspace, and
the matrix $H$ used in Lemma~\ref{lem:cloning} may be reconstructed for the
new pair.  At the end there is a maximizing pair
$(\widehat U,\widehat w)$ whose nonzero-row support $J$ satisfies
\begin{equation}\label{eq:support}
 |J|\le 2^{r-1}d=N_r'.
\end{equation}

A zero row in any maximizing pair has zero weight.  Indeed, its row in the
absolute Gram matrix $B$ is zero, while
\eqref{eq:B-eigenvector} and $\lambda(r,n)\ge1$ imply that the corresponding
weight is zero.  We may consequently delete all indices outside $J$ without
changing feasibility, the weight norm, or the objective.  If $k=|J|$, rank
gives $k\ge r$, and the resulting $k$-row pair gives
\begin{equation}\label{eq:finite-reduction}
 \lambda(r,n)\le\lambda(r,k)\le\lambda(r,N_r'),
\end{equation}
where the last inequality is monotonicity.  Equivalently, one may pad the
$k$-row pair with zero rows and zero weights.

Taking the supremum of \eqref{eq:finite-reduction} over all $n\ge r$ gives
\[
 \lambda(r)\le\lambda(r,N_r')\le\lambda(r),
\]
and hence equality.  Monotonicity then yields
$\lambda(r,n)=\lambda(r)$ for every $n\ge N_r'$.  Finally,
$N_r'\ge r$ for all $r\ge1$; at $r=1$ the argument has one sign cell and
$N_1'=1$.  This proves Theorem~\ref{thm:main}.

\section*{Scope and assistance disclosure}

The theorem improves the explicit universal stability threshold in
\cite{KumarEtAl2026}; it does not determine the least stability index in any
dimension, nor does it address whether that index is polynomial in $r$.
No claim of literature priority beyond the cited comparison is made.

An AI-assisted workflow was used in proof exploration, manuscript drafting,
and consistency checks against the cited maximizing-pair and cloning
arguments.  A separate AI-assisted mathematical review checked the quantified
statement and the preservation, boundary, deletion, and supremum steps; this
does not constitute independent peer review.  No numerical experiment or
computer-generated mathematical certificate is used in the proof.  The author
reviewed the final argument and assumes responsibility for its content.

\bibliographystyle{plain}
\bibliography{references}

\end{document}
